\clearpage\selectlanguage{english}\fancyhead[L]{\small REMARKS.}\pdfbookmark[0]{Remarks}{nachlass-bemerkungen}\SourceLeaf{239}
\begin{center}\textbf{REMARKS.}\par\smallskip\rule{.15\linewidth}{.4pt}\end{center}
In the manuscript, the \textit{Determinatio seriei nostrae per aequationem differentialem etc.} forms the second part of a paper entitled \lq\textit{Disquisitiones generales circa functiones a serie infinita $1+\frac{\alpha.\beta}{1.\gamma}x+$ etc. pendentes auctore} \textsc{Carolo Friderico Gauss} \textit{societati regiae traditae Nov.} 1811.\rq\ After a fuller development of its details, the first part became the well-known paper published in 1812 on January 30. The 27 articles thereby became 37, and the 65 numbered equations were increased to 79. In the manuscript, the numbers of the articles and equations of the second part continue those of the first; for the reader's convenience, however, they have been continued after those of the published paper, and the references that occur have also been related to that paper.

The personal copy of the printed paper \textit{Disqu. gen. circa seriem etc.} contains the following annotation to Article 28:
\[
\begin{aligned}&\textit{The best definition of }\Pi\textit{ is that}\\&
\Pi m=\int_{-\infty}^{+\infty}e^{(m+1)x}e^{-e^x}\,\mathrm dx\end{aligned}
\]
and also equation 58 in this form:
\[
\begin{aligned}&[58]\qquad \log\Pi z=(z+\tfrac12)\log z-z+\tfrac12\log2\pi\\&+\tfrac16.\tfrac12 P+\tfrac18.\tfrac23 Q+\tfrac1{10}.\tfrac34 R\\&+\tfrac1{12}.\tfrac45 S+\text{etc.}\end{aligned}
\]
which can be derived directly from 61, 62 and 57.

In a notebook, with the aid of equation 58 in \textsc{Euler}'s form, the value of $\log\Pi(10+i)$ and, from it, that of $\log\Pi i$ are calculated as
$9{,}7173075-17^\circ16'57''693\,i$ and
$\Pi i=+0{,}4980156-0{,}1549496i$ (after the year 1847).

Article 48 is reproduced here as it must be read in the manuscript after repeated crossings-out of words and entire sentences; according to the author's intention, however, the words \lq\textit{e solo theoremate binomiali}\rq\ should probably still be omitted.
\begin{flushright}\textsc{Schering.}\end{flushright}
\begin{center}\rule{.4\linewidth}{.4pt}\end{center}
